\section{Evaluation}
\label{sec:evaluation}

We evaluate SafeMimic along three research questions:

\begin{itemize}
    \item \textbf{RQ1} (Safety): Does SafeMimic preserve attack functionality across diverse attack chains?
    \item \textbf{RQ2} (Comparison): How does SafeMimic compare against random insertion and regularity-based camouflage (ProvNinja)?
    \item \textbf{RQ3} (Evasion): Does SafeMimic evade both ML-based detectors and rule-based statistical checks?
\end{itemize}

\paragraph{Testbed.}
We deploy a multi-host enterprise network running Ubuntu 22.04 with Linux Audit enabled on bare-metal servers.
Provenance graphs are collected via auditd and reconstructed into CADETS/THEIA/TRACE format.
For detection evasion experiments, we use three DARPA Transparent Computing datasets (CADETS, THEIA, TRACE), each containing system-level audit logs with labeled benign and attack provenance graphs.

\paragraph{Methods under comparison.}
We compare four conditions throughout:
\begin{itemize}
    \item[\textbf{A}] \emph{Raw attack}: unmodified attack chain with no camouflage.
    \item[\textbf{B}] \emph{Random}: camouflage operations drawn uniformly at random from the full benign pool $B_{\text{all}}$, with no safety or distribution constraints.
    \item[\textbf{C}] \emph{ProvNinja}: regularity-based selection~\cite{provninja}, which ranks candidates by frequency in normal behavior. No safety filtering is applied.
    \item[\textbf{D}] \emph{SafeMimic}: our method, selecting from $\Bsafe$ with $\DKL$-minimizing distribution matching and temporal smoothing.
\end{itemize}

%% ============================================================
\subsection{Experiment 1: Attack Preservation (RQ1, RQ2)}
\label{sec:eval:preservation}

\paragraph{Setup.}
We construct four representative APT attack scenarios that span common DARPA TC threat categories:
(1)~\emph{Firefox Backdoor} (CADETS, 5 steps: browser exploit $\to$ payload drop to /tmp $\to$ write to host config $\to$ spawn backdoor process $\to$ establish C2),
(2)~\emph{Phishing APT} (THEIA, 5 steps: credential harvesting (/etc/shadow) $\to$ credential reuse $\to$ network service discovery $\to$ SSH to internal host $\to$ reverse shell),
(3)~\emph{Watering Hole} (TRACE, 4 steps: privilege enumeration (sudo -l) $\to$ SSH key theft $\to$ credential harvesting (/etc/shadow) $\to$ root shell), and
(4)~\emph{Supply Chain} (4 steps: systemd service installation $\to$ library injection $\to$ crontab persistence $\to$ beacon callback).
Each method inserts 5 camouflage operations per attack step.
We repeat every (attack chain, method) pair 100 times with independently sampled camouflage and report the attack success rate---defined as complete execution of all attack steps to their intended final state.

\paragraph{Results.}
\cref{tab:preservation} reports success rates.
SafeMimic achieves 100\% across all four chains, matching the raw attack baseline.
ProvNinja also achieves 100\%, but Random drops to 67\% on Firefox Backdoor and 24\% on Phishing APT.

\begin{table}[t]
\centering
\caption{Attack success rate (\%) over 100 repetitions per cell. SafeMimic preserves attack functionality with a formal guarantee; ProvNinja preserves by coincidence; Random fails on two chains.}
\label{tab:preservation}
\small
\begin{tabular}{lcccc}
\toprule
Method & \rotatebox{55}{Firefox Backdoor} & \rotatebox{55}{Phishing APT} & \rotatebox{55}{Watering Hole} & \rotatebox{55}{Supply Chain} \\
\midrule
A~~Raw Attack    & 100 & 100 & 100 & 100 \\
B~~Random        &  67 &  24 & 100 & 100 \\
C~~ProvNinja     & 100 & 100 & 100 & 100 \\
D~~SafeMimic     & 100 & 100 & 100 & 100 \\
\bottomrule
\end{tabular}
\end{table}

\paragraph{Failure analysis.}
We manually inspect all 109 failures produced by Random insertion (33 in Firefox Backdoor, 76 in Phishing APT).
In Firefox Backdoor, 89\% of failures (29/33) stem from \texttt{kill -9} targeting the compromised process, which terminates the attacker's shell and all child processes.
The remaining 4 failures result from \texttt{systemctl stop \textlangle service\textrangle}, which triggers process termination and has the same effect.
In Phishing APT, 100\% of failures (76/76) trace to \texttt{iptables -A INPUT -j DROP} inserting a deny-all firewall rule on the target host, severing the network path required for reverse-shell establishment.

These failure modes are instructive: the destructive operations are \emph{highly regular} in normal Linux system administration.
Process termination and firewall rule insertion rank in the top 10\% of operations by frequency in our benign workload traces.
A regularity-based selector would rank them highly.
ProvNinja avoids them in these specific chains only because its regularity ranking happens to place other operations (e.g., \texttt{cat /etc/passwd}, \texttt{readdir /proc}) above the destructive ones at the current insertion budget.
This is not a guarantee---increasing the insertion count or changing the benign distribution could cause ProvNinja to select destructive operations.
SafeMimic, by contrast, excludes them \emph{a priori} via $\Seffect(b) \cap \Sdep(A) = \emptyset$ (Theorem~\ref{thm:static_safety}), regardless of the insertion budget or distribution.

%% ============================================================
\subsection{Experiment 2: Detection Evasion (RQ3)}
\label{sec:eval:evasion}

\paragraph{Setup.}
We evaluate evasion on three DARPA TC datasets: CADETS (FreeBSD), THEIA (Linux), and TRACE (Linux).
Each dataset contains 50 benign provenance graphs and 5 attack graphs.
We train a Feature-MLP detector on the benign graphs, extracting 12 structural features per graph (average degree, in/out degree ratio, operation-type entropy, temporal burstiness, etc.) and training a 3-layer MLP (64-32-1) with binary cross-entropy loss.
We also deploy three rule-based checks:
(1)~edge-count threshold ($>\mu + 3\sigma$ of the benign distribution),
(2)~$\DKL$ bound on the operation-type distribution relative to the benign average, and
(3)~a temporal-burst detector that flags windows with $>2\times$ the normal edge rate.
An alert fires if \emph{any} rule triggers.
We apply each evasion method to the 5 attack graphs and report ML detection rate (fraction classified as malicious), average ML anomaly score, rule-alert rate, $\DKL$ from the benign distribution, and edge-count overhead relative to the raw attack graph.

\paragraph{Results.}
\cref{tab:evasion} summarizes the evasion results, averaged across the three datasets.

\begin{table}[t]
\centering
\caption{Detection evasion results averaged across CADETS, THEIA, and TRACE. ML Det: fraction detected by the MLP classifier. ML Score: average anomaly score $\in [0,1]$. Rule Alert: fraction triggering any rule-based check. $\DKL$: KL divergence from benign distribution. Extra Edges: edge-count overhead as a multiple of the raw attack graph.}
\label{tab:evasion}
\small
\begin{tabular}{lccccc}
\toprule
Method & ML Det & ML Score & Rule Alert & $\DKL$ & Extra Edges \\
\midrule
A~~Raw Attack & 100\% & 0.94--1.00 & 80--100\% & 0.15--0.20 & $1\times$ \\
B~~ProvNinja  & 100\% & 1.00       & 100\%     & 0.11--0.12 & $1$--$3\times$ \\
C~~SafeMimic  &   0\% & 0.05       & 80--100\% & 0.06--0.10 & $0.5\times$ \\
\bottomrule
\end{tabular}
\end{table}

\paragraph{ProvNinja backfires.}
The most striking result is that ProvNinja \emph{increases} the ML anomaly score from 0.94--1.00 (raw attack) to a uniform 1.00 across all datasets.
We trace this to three root causes.
First, ProvNinja's regularity-based selection produces a read-heavy operation distribution (${\sim}80\%$ read operations), because reads dominate benign workloads and thus have the highest regularity scores.
This read skew is itself anomalous when injected at the volume ProvNinja requires: real benign behavior has a more balanced read/write/execute mix.
Second, ProvNinja inserts all camouflage operations in a narrow temporal window surrounding the attack steps, producing burst patterns that inflate the temporal-burstiness feature.
Third, many inserted operations target the same small set of resources (e.g., repeatedly reading \texttt{/etc/passwd}), inflating the average in-degree of those nodes and creating a distinct structural signature.
The net effect is that ProvNinja's camouflage creates a \emph{new} anomalous pattern---distinct from both benign behavior and the raw attack---that the ML detector identifies with even higher confidence.

\paragraph{SafeMimic achieves 0\% ML detection.}
SafeMimic's distribution-matched insertion shifts the graph's feature vector toward the benign centroid.
The average $\DKL$ drops to 0.06--0.10, a 55--66\% reduction compared with ProvNinja (0.11--0.12) and a 50--60\% reduction compared with the raw attack (0.15--0.20).
The MLP classifier assigns an average anomaly score of 0.05 (well below the 0.5 decision boundary), yielding a 0\% detection rate.
Temporal smoothing ensures that inserted operations are distributed across the observation window proportionally to the normal edge rate (\cref{fig:temporal}), avoiding burst artifacts.
The operation-type distribution closely tracks the benign profile (\cref{fig:opdist}).

\begin{figure}[t]
    \centering
    \fbox{\parbox{0.9\columnwidth}{\centering\vspace{1.5em}\textbf{[Temporal Edge Rate Plot]}\\[0.3em]
    Baseline (flat) vs ProvNinja (burst) vs SafeMimic (gradual)\vspace{1.5em}}}
    \caption{Temporal edge insertion rate. ProvNinja creates a burst of edges during the evasion window, triggering time-window rule alerts. SafeMimic spreads insertions to match the baseline rate.}
    \label{fig:temporal}
\end{figure}

\begin{figure}[t]
    \centering
    \fbox{\parbox{0.9\columnwidth}{\centering\vspace{1.5em}\textbf{[Operation Distribution Plot]}\\[0.3em]
    Normal vs ProvNinja (read-skewed) vs SafeMimic (matched)\vspace{1.5em}}}
    \caption{Operation type distribution. ProvNinja's camouflage is heavily read-skewed (80\% read operations), creating a distinctive anomalous pattern. SafeMimic matches the baseline distribution.}
    \label{fig:opdist}
\end{figure}

\paragraph{Rule-based alerts are from the attack, not evasion.}
SafeMimic's rule-alert rate (80--100\%) matches the raw attack's, confirming that these alerts originate from the attack operations themselves, not from inserted camouflage.
SafeMimic adds no new rule violations.
ProvNinja, by contrast, pushes the rule-alert rate to 100\% by inflating edge counts ($1$--$3\times$ overhead) and creating temporal bursts.
SafeMimic's edge-count overhead is $0.5\times$---it adds fewer edges than the raw attack contains, because distribution matching favors lightweight operations that each produce few provenance edges.

%% ============================================================
\subsection{Ablation Studies}
\label{sec:eval:ablation}

We conduct four ablation experiments to isolate the contribution of each SafeMimic component.
All ablations use the CADETS dataset with the Firefox Backdoor attack chain unless otherwise noted.

\paragraph{Distribution matching vs.\ random selection from $\Bsafe$.}
Replacing $\DKL$-minimizing selection with uniform random sampling from the safe pool $\Bsafe$ preserves attack functionality (both achieve 100\% success rate, as expected from Theorem~\ref{thm:static_safety}), but increases $\DKL$ from the benign distribution by 40--60\%.
The ML anomaly score rises from 0.05 to 0.31, and detection rate increases from 0\% to 40\%.
This confirms that safety alone is insufficient for evasion; distribution matching is essential.

\paragraph{Static vs.\ dynamic pool.}
We compare SafeMimic's dynamic pool $\Bsafe(t)$, which recomputes safe operations at each attack stage, against a static pool $\Bsafe(t_0)$ fixed at the initial stage.
The dynamic pool contains 15--20\% more operations at later attack stages, because resources consumed during early stages (e.g., temporary files used for exploitation) exit $\Sdep(A)$ once no longer needed.
The larger pool provides more flexibility for distribution matching, reducing $\DKL$ by an additional 8--12\% compared with the static variant.

\paragraph{Number of inserted operations.}
Figure~\ref{fig:nops} shows the effect of varying $n$, the number of camouflage operations inserted per attack step, from 1 to 20.

\begin{figure}[t]
    \centering
    \fbox{\parbox{0.9\columnwidth}{\centering\vspace{1.5em}\textbf{[n\_ops Sensitivity Plot]}\\[0.3em]
    x: n\_ops, y: ML detection rate + KL divergence\vspace{1.5em}}}
    \caption{Effect of the number of inserted operations ($n$). Detection drops to 0\% at $n=5$; beyond $n=10$, diminishing returns with increasing edge-count risk.}
    \label{fig:nops}
\end{figure}
At $n{<}5$, the inserted operations are too few to meaningfully shift the graph's feature distribution, and ML detection remains above 50\%.
Between $n{=}5$ and $n{=}10$, detection drops to 0\% and $\DKL$ reaches its minimum.
Beyond $n{=}10$, additional insertions provide diminishing returns: $\DKL$ plateaus, and the risk of triggering edge-count rules begins to increase.
The sweet spot is $n \in [5, 10]$.

\paragraph{Conservative vs.\ precise effect estimation.}
SafeMimic's default effect estimation is conservative: it over-approximates $\Seffect(b)$ by including all resources of the same type as the operation's target.
A \emph{precise} variant uses resource-name resolution to restrict $\Seffect(b)$ to the exact target resource.
Precise filtering admits 1--3 additional operations per stage---for example, \texttt{kill -9} targeting processes \emph{outside} the attack's process tree, or \texttt{iptables} rules on \emph{other} network interfaces.
These extra operations marginally improve distribution matching ($\DKL$ decreases by 3--5\%), but the conservative variant already achieves 0\% detection.
We default to conservative filtering because it provides a stronger safety margin at negligible cost to evasion effectiveness.

%% ============================================================
\subsection{Cross-Dataset Generalizability}
\label{sec:eval:generalizability}

We evaluate whether SafeMimic's evasion transfers across datasets by training the Feature-MLP detector on one DARPA TC dataset and testing on the remaining two.
\cref{tab:crossdataset} reports the results.

\begin{table}[t]
\centering
\caption{Cross-dataset ML detection rate (\%). Detector trained on the Train set, tested on the Test set. SafeMimic generalizes across datasets; ProvNinja remains fully detected.}
\label{tab:crossdataset}
\small
\begin{tabular}{llcc}
\toprule
Train & Test & ProvNinja & SafeMimic \\
\midrule
CADETS & THEIA & 100 &  0 \\
CADETS & TRACE & 100 & 20 \\
THEIA  & CADETS & 100 &  0 \\
THEIA  & TRACE & 100 & 20 \\
TRACE  & CADETS & 100 &  0 \\
TRACE  & THEIA & 100 &  0 \\
\bottomrule
\end{tabular}
\end{table}

SafeMimic achieves 0--20\% detection across all train/test splits.
The 20\% detection rate on TRACE (when training on CADETS or THEIA) reflects TRACE's distinct system-call vocabulary, which includes FreeBSD-specific operations absent from the Linux-based training data; SafeMimic's distribution matching partially compensates but cannot fully close this gap.
ProvNinja is detected at 100\% in every configuration, confirming that its anomalous insertion pattern is a universal artifact independent of the training distribution.
These results demonstrate that SafeMimic's distribution-matching strategy generalizes: because it targets the \emph{shape} of the benign distribution rather than overfitting to specific features, it remains effective even when the detector is trained on a different dataset.
